\section*{अध्याय \ref{ch:g3:first-electric-circuit} --- पहला विद्युत परिपथ}

\begin{solution}{exo:g3:first-electric-circuit:1}
\omterm{def:g3:first-electric-circuit:battery}{बैटरी} के सिरे: $+$ (धन) और $-$ (ऋण)। \omterm{def:g3:first-electric-circuit:bulb}{बल्ब} के संपर्क: पेंदे के नीचे
की नोक, और पेंदे की धातु वाली बग़ल।
\end{solution}

\begin{solution}{exo:g3:first-electric-circuit:2}
बंद: एक सिरे से दूसरे सिरे तक जाता हुआ एक अटूट फंदा। खुला: वह फंदा
जिसमें कहीं दरार हो। \omterm{def:g3:first-electric-circuit:bulb}{बल्ब} सिर्फ़ \omterm{def:g3:first-electric-circuit:circuit}{बंद परिपथ} में जलता है।
\end{solution}

\begin{solution}{exo:g3:first-electric-circuit:3}
वापसी की यात्रा। एक तार आधा फंदा भर है: बिजली को एक सिरे से निकलकर,
\omterm{def:g3:first-electric-circuit:bulb}{बल्ब} में से होकर, \emph{वापस दूसरे सिरे तक} जाना ही होगा --- लिया को
\omterm{def:g3:first-electric-circuit:bulb}{बल्ब} की बग़ल वाले संपर्क से $-$ सिरे तक एक दूसरा तार चाहिए।
\end{solution}

\begin{solution}{exo:g3:first-electric-circuit:4}
\omterm{ex:g3:first-electric-circuit:switch}{स्विच} फंदे में लगा दरवाज़ा है: बंद हो तो परिपथ पूरा कर देता है, खुला
हो तो दरार बना देता है। बंद \omterm{ex:g3:first-electric-circuit:switch}{स्विच} के भीतर धातु का एक नन्हा पुल दोनों
तारों को जोड़ देता है।
\end{solution}

\begin{solution}{exo:g3:first-electric-circuit:5}
सीधी रेखाओं का एक फंदा, जो इन्हें जोड़ता है: अलग-अलग लंबाई की दो
लकीरें (\omterm{def:g3:first-electric-circuit:battery}{बैटरी}, जिसकी लंबी लकीर $+$ है), बिछा हुआ उठाऊ पुल वाला संकेत
(\omterm{ex:g3:first-electric-circuit:switch}{स्विच}), और क्रॉस वाला एक वृत्त (लैंप)।
\end{solution}

\begin{solution}{exo:g3:first-electric-circuit:6}
चुकी हुई \omterm{def:g3:first-electric-circuit:battery}{बैटरी}; ढीला जोड़ (तार या \omterm{def:g3:first-electric-circuit:bulb}{बल्ब} कसकर न बैठा हो, स्प्रिंग मुड़
गई हो); टूटा हुआ \omterm{def:g3:first-electric-circuit:bulb}{बल्ब} --- उसका दमकता तंतु टूट जाने से ख़ुद \omterm{def:g3:first-electric-circuit:bulb}{बल्ब} के
भीतर ही दरार बन जाती है।
\end{solution}

\begin{solution}{exo:g3:first-electric-circuit:7}
बटन \emph{ही} \omterm{ex:g3:first-electric-circuit:switch}{स्विच} है। उसे दबाने पर धातु का पुल बिछ जाता है और फंदा
बंद हो जाता है, इसलिए घंटी बजती है; छोड़ते ही पुल उठ जाता है, दरार
खुल जाती है, और बजना रुक जाता है।
\end{solution}

\begin{solution}{exo:g3:first-electric-circuit:8}
\omterm{def:g3:first-electric-circuit:battery}{बैटरी} का धक्का बहुत हल्का और सुरक्षित है; दीवार के सॉकेट का धक्का
सैकड़ों गुना ताक़तवर है और चोट पहुँचा सकता है या जान ले सकता है।
सॉकेट प्रयोगों के लिए हैं ही नहीं --- कभी नहीं।
\end{solution}

\begin{solution}{exo:g3:first-electric-circuit:9}
हाँ, वह जलेगा। फंदे का पीछा करो: $+$ सिरा, नोक में से भीतर, दमकते
तंतु से होकर, बग़ल वाले संपर्क से बाहर, तार के सहारे वापस $-$ तक।
यानी दोनों सिरों और दोनों संपर्कों से गुज़रता एक पूरा चक्कर।
\end{solution}

\begin{solution}{exo:g3:first-electric-circuit:10}
उसका फंदा दूसरे सिरे तक जाता ही नहीं। परिपथ को एक सिरे से
\emph{\omterm{def:g3:first-electric-circuit:bulb}{बल्ब} में से होकर दूसरे} सिरे तक जाना चाहिए: $-$ से निकलकर वापस
$-$ में लौटना कहीं न पहुँचने वाला चक्कर है --- \omterm{def:g3:first-electric-circuit:battery}{बैटरी} के धक्के को अपने
दोनों अलग सिरे खेल में चाहिए।
\end{solution}

\begin{solution}{exo:g3:first-electric-circuit:11}
सारे \omterm{def:g3:first-electric-circuit:bulb}{बल्ब} एक ही अकेले फंदे पर, एक के बाद एक लगे होंगे: ग़ायब \omterm{def:g3:first-electric-circuit:bulb}{बल्ब} एक
दरार छोड़ जाता है, और एक दरार पूरी लड़ी को अँधेरा कर देती है ---
यही फंदे का नियम है। (अगले साल इस एक-के-बाद-एक वाली तरतीब को उसका
नाम मिलेगा: श्रेणी परिपथ।)
\end{solution}
